%
% $Log: conclusion.tex,v $
% Revision 1.3  2002/01/08 18:52:29  mjm
% that
%
% Revision 1.2  1999/07/28 17:23:11  mjm
% version sent back to Kristian
%
% Revision 1.1  1999/06/10 23:33:32  mjm
% Initial revision
%
%

\handout{Conclusions}

Your solution should end with a brief concluding paragraph. The
conclusion is a short, concise paragraph that summarizes the main
ideas, techniques and results discussed in the problem.  The
concluding paragraph should focus on the main results and main ideas
so that the reader can absorb the important parts of your solution to
the problem. The concluding paragraph should not
introduce any new material but simply bring back the reader to the main ideas
of your solution to the problem. This means that you should not
discuss something in the conclusion that was not discussed earlier in
your solution.

When you write a conclusion, put yourself in the potential reader's
position and imagine what kind of information this reader may look for
when he or she reads the concluding paragraph.  Also, ask yourself
``What do I want my reader to remember?''  Here are a few suggestions
of what you may want to include in the conclusion paragraph:

\begin{itemize}
\item Restate the problem you solved.
\item Restate the results and interpretation of the results. If your
  results are lengthy, such as a big table with numerical values,
  just  describe the results qualitatively. For example,
  give the range of values for your results.
\item Indicate if your results seem reasonable. If not, then explain why.
\item If you were unsuccessful in solving your problem, mention
  possible reasons for this.
\item If applicable, mention other problems that are related to the
  problem you solved.
\item Give suggestions for improvements, i.e., what else you could do
  for the problem.
\end{itemize}

These are just suggestions, and not all of them may be applicable for
your problem.  As mentioned above, the conclusion should be short and
concise without discussing anything new that was not discussed earlier
in the report.


\secondpage

\subsection*{Examples:}

{\bf Problem:}
  \emph{Find
        the dimensions of a right circular cylinder of volume
        $10^{-6}\ m^3$ such that a minimum amount of material is
        used.}  

\bigskip

      \begin{badexample}
        In this problem I used calculus to solve a mathematical
        problem.\\
\fbox{This gives no information.}
    \end{badexample}

\bigskip

    \begin{badexample}
      I solved for $h$ in $\pi r^2h=10^{-6}$ and plugged the
      expression for $h$ into $2\pi r^2+2\pi rh$. Then I took the
      derivative of $2\pi r^2+\frac{2\times 10^{-6}}{r}$. Then I set
      the derivative equal to zero. Then I solved for $r$.
    \end{badexample}


\noindent\fbox{\parbox{\textwidth}{\rm 
      The last example fails to focus on the main ideas of the
      problem. It gives a (bad) summary of some of the steps used to
      solve the problem. It does not summarize the problem given to us
      and it does not clearly indicates what your bottom line was or
      if your were even successful in solving the problem.
}}

\bigskip

    \begin{goodexample}
      The problem discussed in this paper is to find the dimensions of
      a right circular cylinder of volume $10^{-6}\ m^3$ such that a
      minimum amount of material is used. The problem was solved by
      minimizing the area of the cylinder using the first and second
      derivative test. The optimal dimensions were found to be radius
      $\cong 5.42\ cm$ and height $\cong 10.84\ cm$. The method used
      for solving the problem can be applied to other optimization
      problems, e.g., finding the dimensions of other shapes with a
      fixed volume such that the used material is minimized. A
      possible development of the technique used could be to consider
      optimization of a function of more than two parameters and more
      than one constraint.
    \end{goodexample}

\noindent\fbox{\parbox{\textwidth}{\rm 
      Note that the suggested conclusion requires that we actually did
      discuss possible applications and developments of the technique
      earlier in the report, something that may not be required for
      you in your course.
}}

\marginpar{conclusion}